This paper introduces a novel, mechanistically grounded framework for using acoustic and vibrotactile stimuli as direct interventions in human physiology: \textbf{Resonant Bioharmonic Medicine (RBM)}. Drawing from recent advances in biophysics, neuroscience, and cellular biomechanics, I argue that sound can be used not merely for emotional modulation, but as a \emph{biological compiler}---a tool to entrain, amplify, or re-harmonize specific metabolic and neural processes at their natural frequencies.

The foundation of RBM is a layered model of the human body as a system of coupled oscillators, extending from intracellular organelles to neural field dynamics. Within this framework, published studies linking specific frequency bands---including 40\,Hz, 0.8\,Hz, and 100\,Hz---to measurable physiological outcomes are especially significant. Reported effects include reduced tau accumulation, improved heart-rate variability (HRV), enhanced sleep-phase memory consolidation, and cytoskeletal modulation.

Building on this evidence, I propose the design of localized ``frequency medicine'' devices that apply phase-locked, symbolic, or biologically matched vibratory stimuli directly to the body. Multiple low-risk, high-impact prototypes can be constructed immediately using off-the-shelf components, including vagal nerve acoustic entrainment, sleep-aligned audio fields, and mitochondrially resonant stimulation pads.

This paper concludes by framing RBM as a bridge between bioelectromagnetic physiology, symbolic cognition, and intentional systems design, with potential implications for neurodegenerative disease, trauma recovery, and perceptual engineering.
